\section{Scope, Terms, Evidence, and Limits}

\subsection{The object and the reader}

At its narrowest, \emph{Septuagint} names the Greek translation of the Torah---Genesis through Deuteronomy---whose origin was remembered in the \emph{Letter of Aristeas}. In textual scholarship, \emph{Old Greek} usually means the earliest recoverable Greek translation of a particular Hebrew or Aramaic book, whether or not it belonged to that first project. In ordinary Christian usage, \emph{Septuagint} can mean the wider Greek Old Testament transmitted by the Church: translations, books composed in Greek, expanded Greek forms, and later revisions gathered in different combinations. Finally, a printed “Septuagint” is an editor's text reconstructed from manuscripts that are much later than the first translations. These meanings overlap; they are not interchangeable.

This account asks four connected questions. What can be recovered about the production of the Greek Torah? How did further scriptural books enter Greek? Why do their texts differ from later Hebrew forms and from one another? How did Jewish translations become Christian Scripture without ceasing to be evidence for ancient Jewish interpretation?

It is written for a serious reader who may know the Septuagint chiefly through claims about the “Bible of Jesus,” the deuterocanonical books, or differences from the Masoretic Text. It does not adjudicate inspiration or define a canon. It reconstructs the histories that those theological questions presuppose.

The historical range begins with the Greek Torah in Ptolemaic Egypt in the third century BC and follows later translation, Jewish revision, Christian codices and reception, and modern critical reconstruction through the publication cutoff. Its geographic center is Egypt and the eastern Mediterranean; other regions enter only where a manuscript, ancient witness, daughter version, or modern edition materially bears on that history. Daughter versions are used as textual witnesses, not given complete independent histories.

\subsection{Terms that prevent false precision}

\begin{threestable}{Term}{Meaning in this account}{Boundary}
LXX / Septuagint & The Greek Torah when discussing origins; the wider Greek scriptural corpus when context makes that usage explicit. & Never a claim that all books share one translator, date, source text, or textual history.\\
Old Greek (OG) & The earliest Greek form recoverable for a particular translation unit. & A scholarly reconstruction, not necessarily an extant manuscript or a claim of verbal access to the translator's copy.\\
Source text / \emph{Vorlage} & The Hebrew or Aramaic form from which a Greek translator worked. & Not automatically identical with the medieval Masoretic Text or any surviving scroll.\\
Masoretic Text (MT) & The consonantal tradition with medieval vocalization and notes represented in standard Hebrew editions. & A central witness to the Hebrew Bible, not a time machine into every translator's exemplar.\\
Revision / recension & Deliberate alteration of an existing Greek version, often toward a Hebrew reference text or an editorial standard. & Must be demonstrated by patterns; not every variant is a recension.\\
Great codex & A large Christian parchment Bible such as Vaticanus, Sinaiticus, or Alexandrinus. & A fourth- or fifth-century collection, not a direct inventory of third-century Alexandrian Judaism.\\
\end{threestable}

\subsection{Method and limits}

The earliest physical witnesses are fragmentary and the full codices are Christian and centuries later. Ancient authors write for apologetic, historical, or theological purposes. Modern editors must reason backward through revision and scribal transmission. Accordingly, this study distinguishes four evidence classes: early artifacts; ancient narrative witnesses; later Jewish and Christian reception; and modern text-critical reconstruction. A date stated as a range marks the limits of the evidence rather than indecision to be concealed.

No manuscript image, facsimile, modern translation, or extended third-party passage is incorporated. Brief ancient wording, where needed, is attributed to an identified public-domain human translation; modern translations are paraphrased. Online records and mutable project descriptions were checked through 15 July 2026. The source audit and evidence map record the full claim-to-source alignment, exact loci, access dates, negative results, and rights decisions.

\claimcheck{“The Septuagint says”}{A responsible use of this phrase must identify at least the book and, when the reading matters, the Greek form or edition. Greek Daniel is not one simple text; Greek Tobit is pluriform; the Old Greek of Job was supplemented; and the large codices disagree. The singular title is convenient only after the historical plurality is made visible.}

\subsection{Publication boundary}

This is a source-audited historical study, checked through 15 July 2026. It is not a critical edition, a project translation of Scripture, or an ecclesiastical judgment about canon or inspiration. Independent historical and text-critical review remains outstanding.
